\documentclass{article}
\usepackage[utf8]{inputenc}
\usepackage{amsmath}
\usepackage{amsthm}
\usepackage{amssymb,amsfonts}
\usepackage{mathrsfs}
\usepackage[hidelinks]{hyperref}
\usepackage[capitalize]{cleveref}
\usepackage{stmaryrd}

\newtheorem{theorem}{Theorem}[section]
\newtheorem{proposition}[theorem]{Proposition}
\newtheorem{claim}[theorem]{Claim}
\newtheorem{corollary}[theorem]{Corollary}
\newtheorem{lemma}[theorem]{Lemma}
\theoremstyle{definition}
\newtheorem{definition}[theorem]{Definition}
\newtheorem{convention}[theorem]{Convention}
\newtheorem{remark}[theorem]{Remark}

\newcommand{\Rbar}{\overline{\mathbb{R}}}
\newcommand{\R}{\mathbb{R}}
\newcommand{\rk}{\operatorname{rk}}
\newcommand{\cl}{\operatorname{cl}}
\newcommand{\Trop}{\operatorname{Trop}}
\newcommand{\op}{\operatorname{op}}
\newcommand{\into}{\hookrightarrow}
\newcommand{\onto}{\twoheadrightarrow}
\newcommand{\Dr}{\mathbf{Dr}}
\newcommand{\bV}{\mathbf V}
\newcommand{\cLL}{\mathcal{L}}
\newcommand{\cW}{\mathcal{W}}
\newcommand{\PT}[1]{\mathbf{P}\big(\Rbar^{#1}\big)}

\title{First proof: Duality for tropical linear subspaces}
\author{}
\date{}

\begin{document}
\maketitle

\noindent
Let $M$ be a matroid on the ground set $E=[n]$, of rank $r$, and let $\Trop(M)$ be the
corresponding tropical linear space, i.e.\ the closure of the Bergman fan of $M$ in tropical
projective space $\PT{n}$. The \emph{relative Dressian} of $M$ is the set
\[
  \Dr(M):=\{\text{tropical linear subspaces of }\Trop(M)\},
\]
partially ordered by inclusion of tropical linear spaces. Suppose $M$ admits an order-reversing
involution $F\mapsto F^{\perp}$ on its lattice of flats $\cLL(M)$. There is a natural inclusion of
posets $\cLL(M)^{\op}\into\Dr(M)$ given by $F\mapsto\Trop(M/F\oplus U_{0,F})$.

\medskip
\noindent\textbf{Theorem.} \emph{There is an order-reversing involution $\Phi$ on $\Dr(M)$ that
extends the involution $F\mapsto F^{\perp}$; that is, $\Phi$ restricts to $F\mapsto F^{\perp}$ on the
image of $\cLL(M)^{\op}\into\Dr(M)$.}

\medskip

We first fix conventions and isolate the two standard foundational facts we rely on
(\Cref{ssec:found}). We then prove the rank identity forced by $\perp$ (\Cref{lem:rank}); show that
every element of $\Dr(M)$ has a canonical \emph{flat-coordinate} description, in that the underlying
valuated quotient is constant on the bases lying in a fixed flat (\Cref{prop:constancy}); make this
description into a \emph{cryptomorphism} by writing down the explicit \emph{essential relations} that
characterise flat-coordinates of quotients and proving that they are equivalent to the full system
(\Cref{prop:essential-relations}); define $\Phi$ explicitly by precomposing flat-coordinates with
$\perp$ (\Cref{def:Phi}); and verify that $\Phi$ is well-defined (\Cref{prop:wd}), an involution
(\Cref{prop:inv}), order-reversing (\Cref{prop:rev}), and extends $\perp$ (\Cref{prop:ext}). The
honest status of each ingredient is summarised in \Cref{rmk:honest}.

\section{Conventions and foundational facts}\label{ssec:found}

We use the min-plus conventions on $\Rbar=\R\cup\{\infty\}$. A \emph{valuated matroid} of rank $k$
on $E$ is a function $\mu\colon\binom{E}{k}\to\Rbar$, not identically $\infty$, satisfying the
tropical Pl\"ucker relations: for every $S\in\binom{E}{k-1}$ and every $T\in\binom{E}{k+1}$,
\begin{equation}\label{eqn:plucker}
  \min_{a\in T\setminus S}\bigl\{\mu(S+a)+\mu(T-a)\bigr\}\ \text{vanishes tropically,}
\end{equation}
i.e.\ either all summands are $\infty$, or the minimum is attained for at least two
$a\in T\setminus S$. Here $S+a:=S\cup\{a\}$ and $T-a:=T\setminus\{a\}$. Two valuated matroids of the
same rank are \emph{equivalent} if they differ by a global additive constant; the support
$\{B:\mu(B)<\infty\}$ is the set of bases of the \emph{underlying matroid} $\underline\mu$. The
\emph{trivial valuation} $\underline M$ of $M$ is the rank-$r$ valuated matroid with $\underline
M(B)=0$ if $B$ is a basis of $M$ and $\underline M(B)=\infty$ otherwise; thus
$\underline{\underline M}=M$.

\begin{definition}\label{def:quotient}
Let $\nu,\mu$ be valuated matroids on $E$ of ranks $k\le m$. We say $\nu$ is a \emph{valuated
quotient} of $\mu$, written $\mu\onto\nu$, if for every $S\in\binom{E}{k-1}$ and every
$T\in\binom{E}{m+1}$,
\begin{equation}\label{eqn:incidence-plucker}
  \min_{a\in T\setminus S}\bigl\{\nu(S+a)+\mu(T-a)\bigr\}\ \text{vanishes tropically.}
\end{equation}
Here $\nu$ is evaluated on the $k$-set $S+a$ and $\mu$ on the $m$-set $T-a$; the index sets have the
correct cardinalities $|S+a|=k$ and $|T-a|=m$, since $|S|=k-1$ and $|T|=m+1$. When $m=k+1$ this is
the \emph{elementary} quotient relation; the relation \eqref{eqn:incidence-plucker} for general
$k\le m$ is equivalent to the existence of a chain of elementary quotients
$\mu=\mu_m\onto\mu_{m-1}\onto\cdots\onto\mu_k=\nu$ (Murota, Brandt--Eur--Speyer--Sturmfels).
\end{definition}

\begin{proposition}[Dressian as poset of valuated quotients]\label{fact:dressian}
Let $M$ be loopless of rank $r$. The assignment $\theta\mapsto\Trop\theta$ is an isomorphism of
posets
\[
  \bV(M)/\!\sim\ \xrightarrow{\ \cong\ }\ \Dr(M),
\]
where $\bV(M):=\bigsqcup_{0\le k\le r}\bV^{k}(M)$, $\bV^{k}(M)$ is the set of corank-$k$ valuated
quotients of $\underline M$ (i.e.\ valuated quotients of rank $r-k$), $\sim$ is equivalence up to a
global additive constant, the order on $\bV(M)/\!\sim$ is the quotient relation
($\theta_2\ge\theta_1\iff\theta_2\onto\theta_1$), and the order on $\Dr(M)$ is inclusion. In
particular $\theta_2\onto\theta_1\iff\Trop\theta_1\subseteq\Trop\theta_2$.
\end{proposition}

\begin{proof}
This is the relative form of the standard correspondence between tropical linear subspaces and
valuated matroids, and between inclusions and valuated quotients (Maclagan--Sturmfels,
\emph{Introduction to Tropical Geometry}, Ch.~4; Brandt--Eur--Sturmfels; Brandt--Speyer). We verify
its standing hypothesis that $M$ is loopless. A loop of $M$ lies in $\cl(\varnothing)$, the bottom
flat $\hat0$ of $\cLL(M)$. The order-reversing bijection $\perp$ sends the unique top
$\hat1=E$ to the unique bottom, so $\hat1^{\perp}=\hat0$ and $\hat0^{\perp}=\hat1$; in particular
$\cLL(M)$ has length $r$ from $\hat0$ to $\hat1$, so $\hat0$ has rank $0$ and contains no nonloop.
If $M$ had loops $L$, then $\Trop(M)$, $\cLL(M)$, $\perp$, and $\Dr(M)$ are unchanged upon deleting
$L$ (loops are common to all bases' complements and carry no coordinate in $\Trop M$); we may thus
assume $M$ loopless. The order-preservation in both directions is the cited quotient/containment
dictionary.
\end{proof}

\begin{proposition}[Valuated matroid duality]\label{fact:dual}
Let $\theta$ be a valuated matroid of rank $k$ on $E$, $|E|=n$. Define
$\theta^{*}\colon\binom{E}{n-k}\to\Rbar$ by $\theta^{*}(B):=\theta(E\setminus B)$. Then $\theta^{*}$
is a valuated matroid of rank $n-k$ with $\underline{\theta^{*}}=(\underline\theta)^{*}$ and
$(\theta^{*})^{*}=\theta$, and $\mu\onto\nu$ implies $\nu^{*}\onto\mu^{*}$.
\end{proposition}

\begin{proof}
Dress--Wenzel duality. The complementation $B\mapsto E\setminus B$ identifies $\binom{E}{n-k}$ with
$\binom{E}{k}$ and, for $a\in T\setminus S$, satisfies $E\setminus(S+a)=(E\setminus S)-a$ and
$E\setminus(T-a)=(E\setminus T)+a$; thus it carries the Pl\"ucker relation \eqref{eqn:plucker} of
$\theta$ at $(S,T)$ term-by-term to the Pl\"ucker relation of $\theta^{*}$ at $(E\setminus
T,E\setminus S)$, and similarly carries the incidence relation \eqref{eqn:incidence-plucker} for
$\mu\onto\nu$ to that for $\nu^{*}\onto\mu^{*}$. Double duality is immediate from $E\setminus
(E\setminus B)=B$.
\end{proof}

\begin{convention}\label{conv:perp}
Throughout, $\perp$ is a fixed inclusion-reversing bijection $\cLL(M)\to\cLL(M)$ with
$(F^{\perp})^{\perp}=F$. We do \emph{not} assume $M$ is a modular matroid: classical modularity of a
matroid (every pair of flats is a modular pair) is a strictly stronger condition than the existence
of $\perp$, and we use only the existence of $\perp$ and its consequences derived below.
\end{convention}

\section{The rank identity}

\begin{lemma}\label{lem:rank}
For every $F\in\cLL(M)$, $\rk_M(F)+\rk_M(F^{\perp})=r$. Moreover $\hat0^{\perp}=\hat1$ and
$\hat1^{\perp}=\hat0$, where $\hat0=\cl(\varnothing)$, $\hat1=E$. The map $\perp$ is an
anti-automorphism of the lattice $\cLL(M)$: $(F\vee G)^{\perp}=F^{\perp}\wedge G^{\perp}$ and
$(F\wedge G)^{\perp}=F^{\perp}\vee G^{\perp}$.
\end{lemma}

\begin{proof}
$\cLL(M)$ is graded of rank $r$: every maximal chain $\hat0=G_0\lessdot\cdots\lessdot G_r=\hat1$ has
length $r$, and $\rk_M(G)$ equals the height of $G$. An order-reversing bijection sends covers to
covers ($A\lessdot B\iff B^{\perp}\lessdot A^{\perp}$), hence maximal chains to maximal chains in
reverse order. Applying $\perp$ to a maximal chain through $F$ at height $\rk_M(F)$ yields a maximal
chain in which $F^{\perp}$ sits at height $r-\rk_M(F)$, so $\rk_M(F^{\perp})=r-\rk_M(F)$. Taking
$F=\hat0$ gives $\hat0^{\perp}=\hat1$, and involutivity gives $\hat1^{\perp}=\hat0$. Finally, in any
lattice an order-reversing bijection exchanges joins and meets, because $F\vee G$ is the least upper
bound of $\{F,G\}$, whose image under the order-reversing $\perp$ is the greatest lower bound of
$\{F^{\perp},G^{\perp}\}$, namely $F^{\perp}\wedge G^{\perp}$; and dually.
\end{proof}

\section{Flat-coordinates: constancy on flats}

Write $\cLL_k(M)$ for the rank-$k$ flats of $M$.

\begin{proposition}[Constancy on flats]\label{prop:constancy}
Let $\theta\in\bV^{\,r-k}(M)$ be a valuated quotient of $\underline M$ of rank $k$, with underlying
matroid $N:=\underline\theta$. If $B,B'$ are bases of $N$ with $\cl_M(B)=\cl_M(B')$, then
$\theta(B)=\theta(B')$. Consequently the function
\begin{equation}\label{eqn:hat-def}
  \widehat\theta\colon\cLL_k(M)\to\Rbar,\qquad
  \widehat\theta(F):=\begin{cases}\theta(B), & \text{$F=\cl_M(B)$ for a basis $B$ of $N$},\\
  \infty, & \text{no basis $B$ of $N$ has $\cl_M(B)=F$,}\end{cases}
\end{equation}
is well-defined up to the global additive constant of $\theta$, and $\theta(B)=\widehat\theta(\cl_M
B)$ for every basis $B$ of $N$. (If $B$ is a basis of $N$ then $B$ is independent in $M$, so
$\cl_M(B)\in\cLL_k(M)$.)
\end{proposition}

\begin{proof}
Since $M\onto N$ is a matroid quotient, $\rk_N\le\rk_M$ and every $N$-independent set is
$M$-independent; in particular a basis $B$ of $N$ (size $k$) is $M$-independent, so
$\cl_M(B)\in\cLL_k(M)$. This proves the parenthetical claim.

For the constancy, it suffices to treat $B'=B-x+y$ obtained from $B$ by a single symmetric exchange
($x\in B$, $y\notin B$, $B'$ a basis of $N$): the basis-exchange graph of $N$ is connected, and we
may restrict exchanges to stay inside the rank-$k$ flat $F:=\cl_M(B)=\cl_M(B')$, since $B,B'$ are
both bases of the restriction $N|_F$ and the basis-exchange graph of the matroid $N|_F$ is connected
with all bases inside $F$. So assume $B'=B-x+y$ with $\cl_M(B)=\cl_M(B')=F$; thus
$y\in\cl_M(B)\setminus B$, i.e.\ $B+y=B'\cup\{x\}$ has $M$-rank $k$ (it lies in $F$).

Apply the quotient relation \eqref{eqn:incidence-plucker} for $\underline M\onto\theta$ (so $m=r$,
the relation of $\theta$ with the trivial valuation) with
\[
  S:=B-x\quad(|S|=k-1),\qquad T:=(B-x)\cup\{x,y,z_1,\dots,z_{r-k}\}\quad(|T|=r+1),
\]
where $z_1,\dots,z_{r-k}\in E$ are chosen so that $C:=\{x,z_1,\dots,z_{r-k}\}$ extends $B$ to a basis
of $M$, i.e.\ $S\cup C=B\cup\{z_1,\dots,z_{r-k}\}$ is a basis of $M$. Such $z_i$ exist because
$\rk_M(B)=k$ and $M$ has rank $r$ (extend the independent set $B$ to an $M$-basis). Note $z_i\notin
F$ for all $i$, since adding any $z_i$ raises $M$-rank above $k$.

The relation \eqref{eqn:incidence-plucker} states that
\[
  \min_{a\in T\setminus S}\bigl\{\theta(S+a)+\underline M(T-a)\bigr\}\ \text{vanishes tropically,}
  \qquad T\setminus S=\{x,y,z_1,\dots,z_{r-k}\}.
\]
We compute the summands $\underline M(T-a)$, which is $0$ if $T-a$ is a basis of $M$ and $\infty$
otherwise. The set $T-a$ has $r$ elements.
\begin{itemize}
\item For $a=y$: $T-y=(B-x)\cup\{x,z_1,\dots,z_{r-k}\}=B\cup\{z_1,\dots,z_{r-k}\}$, a basis of $M$
by choice; so $\underline M(T-y)=0$, and $\theta(S+y)=\theta(B-x+y)=\theta(B')$.
\item For $a=x$: $T-x=(B-x)\cup\{y,z_1,\dots,z_{r-k}\}=(B'-y)\cup\{y,z_1,\dots,z_{r-k}\}$; since
$B'=B-x+y$ is $M$-independent of rank $k$ and $\{z_1,\dots,z_{r-k}\}$ extends $B$ (equivalently
$F=\cl_M B=\cl_M B'$) to an $M$-basis, $T-x=B'\cup\{z_1,\dots,z_{r-k}\}$ is also a basis of $M$
(it spans, has $r$ elements, and is independent because $B'$ spans $F$ and the $z_i$ are independent
modulo $F$). So $\underline M(T-x)=0$ and $\theta(S+x)=\theta(B)$.
\item For $a=z_j$: $T-z_j=(B-x)\cup\{x,y,z_1,\dots,\widehat{z_j},\dots,z_{r-k}\}=B\cup\{y\}\cup
\{z_i:i\ne j\}$. Since $y\in F=\cl_M(B)$, this set has $M$-rank $\le k+(r-k-1)=r-1<r$, hence is
\emph{not} a basis of $M$; so $\underline M(T-z_j)=\infty$, and the summand for $a=z_j$ is $\infty$.
\end{itemize}
Thus the only finite summands are those for $a=x$ and $a=y$, equal to $\theta(B)$ and $\theta(B')$
respectively. For the tropical minimum to vanish (be attained at least twice, or be $\infty$), we
need $\theta(B)=\theta(B')$ (if either is finite; if both are $\infty$ the claim is trivial). Hence
$\theta(B)=\theta(B')$, completing the proof of constancy. Well-definedness of $\widehat\theta$ and
the identity $\theta(B)=\widehat\theta(\cl_M B)$ follow immediately.
\end{proof}

\Cref{prop:constancy} attaches to each $\theta\in\bV^{\,r-k}(M)$ its flat-coordinate
$\widehat\theta\in\cW_k(M)$, where $\cW_k(M)$ denotes the set of functions
$\cLL_k(M)\to\Rbar$ (not identically $\infty$) modulo global additive constant. Conversely, to a
function $w\colon\cLL_k(M)\to\Rbar$ we associate the function
\begin{equation}\label{eqn:thetaw}
  \theta_w\colon\binom{E}{k}\to\Rbar,\qquad
  \theta_w(B):=\begin{cases}w(\cl_M B), & \rk_M(B)=k,\\ \infty, & \rk_M(B)<k.\end{cases}
\end{equation}
A function $\theta\colon\binom{E}{k}\to\Rbar$ is called \emph{flat-constant} (with respect to $M$) if
$\theta=\theta_w$ for some $w$; equivalently, $\theta(B)<\infty\Rightarrow\rk_M(B)=k$ and $\theta$
is constant on the fibres of $B\mapsto\cl_M(B)$. By \Cref{prop:constancy} every
$\theta\in\bV^{\,r-k}(M)$ is flat-constant and equals $\theta_{\widehat\theta}$ on its support; and
$w\mapsto\theta_w$, $\theta\mapsto\widehat\theta$ are mutually inverse between flat-constant
functions and functions on $\cLL_k(M)$ (up to the support convention $\widehat{\theta_w}=w$ on
$\operatorname{supp}w$). We write
\[
  \cW(M):=\bigsqcup_{0\le k\le r}\cW_k(M).
\]

\section{The cryptomorphism: essential relations}\label{sec:crypto}

We now make the flat-coordinate description into a cryptomorphism: we write down explicitly which
functions $w\in\cW_k(M)$ are the flat-coordinates of valuated quotients of $\underline M$, in a form
that is manifestly invariant under the lattice anti-automorphism $\perp$. This is the proposition on
which the entire construction rests; we state the essential relations, prove they cut out exactly the
flat-coordinates of quotients, and record their $\perp$-invariance.

Recall that an interval $[A,Z]$ of $\cLL(M)$ with $\rk_M(A)=k-1$, $\rk_M(Z)=k+1$ has, by gradedness,
an associated set of \emph{covering atoms}
\[
  \mathrm{At}(A,Z):=\{\,F\in\cLL_k(M):A\lessdot F\lessdot Z\,\}
  =\{\,F\in\cLL_k(M):A<F<Z\,\},
\]
the rank-$k$ flats strictly between $A$ and $Z$. (In matroid terms $\mathrm{At}(A,Z)$ is the set of
lines of the rank-$2$ interval $[A,Z]$ of the geometric lattice $\cLL(M)$.)

\begin{definition}[Essential relations]\label{def:essential}
A function $w\colon\cLL_k(M)\to\Rbar$, not identically $\infty$, satisfies the \emph{essential
relations of rank $k$} if:
\begin{enumerate}
\item[(E1)] (\emph{support is a quotient}) $\operatorname{supp}w:=\{F\in\cLL_k(M):w(F)<\infty\}$ is
the set of rank-$k$ flats of a matroid quotient $N$ of $M$ on $E$ (equivalently, $\theta_w$ has
nonempty support and $\underline{\theta_w}$ is a quotient of $M$ in the ordinary matroid sense);
\item[(E2)] (\emph{tropical Pl\"ucker on modular intervals}) for every interval $[A,Z]$ of
$\cLL(M)$ with $\rk_M(A)=k-1$ and $\rk_M(Z)=k+1$,
\begin{equation}\label{eqn:essential}
  \min_{F\in\mathrm{At}(A,Z)}\,w(F)\ \text{vanishes tropically,}
\end{equation}
i.e.\ either $w(F)=\infty$ for all $F\in\mathrm{At}(A,Z)$, or the minimum is attained for at least
two $F\in\mathrm{At}(A,Z)$.
\end{enumerate}
We write $\cW^{\mathrm{ess}}_k(M)$ for the set of such $w$, modulo global additive constant, and
$\cW^{\mathrm{ess}}(M):=\bigsqcup_k\cW^{\mathrm{ess}}_k(M)$.
\end{definition}

The relations (E2) are exactly the three-term tropical Pl\"ucker relations on the geometric lattice
$\cLL(M)$ attached to modular pairs $(A,Z)$ of ranks $k-1$ and $k+1$. There is one for each such
interval, so the list is finite once $M$ is fixed.

\begin{proposition}[Cryptomorphism]\label{prop:essential-relations}
Let $M$ be loopless of rank $r$ and $0\le k\le r$. The map $\theta\mapsto\widehat\theta$ is a
bijection
\[
  \bV^{\,r-k}(M)/\!\sim\ \xrightarrow{\ \cong\ }\ \cW^{\mathrm{ess}}_k(M),
\]
with inverse $w\mapsto\theta_w$. Equivalently: a flat-constant function $\theta=\theta_w$ is a
valuated quotient of $\underline M$ of rank $k$ \emph{if and only if} $w$ satisfies the essential
relations \textup{(E1)--(E2)}. In particular $\cW^{\Dr}(M)=\cW^{\mathrm{ess}}(M)$ inside $\cW(M)$,
and
\begin{equation}\label{eqn:identification}
  \Dr(M)\ \cong\ \bV(M)/\!\sim\ \cong\ \cW^{\mathrm{ess}}(M)\quad\text{as posets,}
\end{equation}
the last identification carrying the quotient order on $\bV(M)/\!\sim$.
\end{proposition}

The proof occupies the rest of this section. It has two parts: the essential relations are
\emph{necessary} (every quotient gives such a $w$), and \emph{sufficient} (they imply the full
tropical Pl\"ucker and incidence systems). The sufficiency is exactly the assertion that all
non-essential relations are redundant.

\begin{lemma}[(E2) is the incidence relation in flat coordinates]\label{lem:E2-incidence}
Let $w\colon\cLL_k(M)\to\Rbar$ and $\theta_w$ as in \eqref{eqn:thetaw}. Then $w$ satisfies
\textup{(E2)} if and only if $\theta_w$ satisfies the elementary incidence relations with
$\underline M$, namely \eqref{eqn:incidence-plucker} with $\nu=\theta_w$, $\mu=\underline M$,
$m=r$:
\begin{equation}\label{eqn:incid-r}
  \min_{a\in T\setminus S}\bigl\{\theta_w(S+a)+\underline M(T-a)\bigr\}\ \text{vanishes tropically}
  \qquad\bigl(S\in\tbinom{E}{k-1},\ T\in\tbinom{E}{r+1}\bigr).
\end{equation}
\end{lemma}

\begin{proof}
Fix $S\in\binom{E}{k-1}$, $T\in\binom{E}{r+1}$. In \eqref{eqn:incid-r} the summand for $a$ is finite
only if both $\theta_w(S+a)<\infty$ and $\underline M(T-a)=0$, i.e.\ $S+a$ is $M$-independent of rank
$k$ \emph{and} $T-a$ is a basis of $M$. We claim the set of such $a$, and the associated values
$w(\cl_M(S+a))$, are governed precisely by an interval $[A,Z]$ as in (E2), and conversely every such
interval arises this way; the two minima then coincide term-by-term, proving the equivalence.

\emph{From $(S,T)$ to an interval.} Suppose some summand of \eqref{eqn:incid-r} is finite, witnessed
by $a_0$. Then $T-a_0$ is a basis of $M$ and $S+a_0$ is independent. We may assume $S$ is
$M$-independent (otherwise $\rk_M(S+a)\le\rk_M(S)<k$ for all $a\notin S$ forces every $\theta_w$-term
to be $\infty$, and likewise on the (E2) side $\mathrm{At}(A,Z)$ is empty when $\rk_M(A)<k-1$; both
minima are then ``all $\infty$'' and vacuously vanish). Set $A:=\cl_M(S)$, of rank $k-1$. For the
summand of $a$ to be finite we need $\underline M(T-a)=0$, i.e.\ $T-a$ a basis; since $|T|=r+1$,
exactly the $a\in T$ with $T-a$ a basis are those for which $a$ is not a coloop of $M|_T$, and these
$a$ together with $A$ generate, via $\cl_M(S+a)$, the rank-$k$ flats covering $A$ inside the rank of
$T$. Concretely, the elements $a\in T\setminus S$ with $T-a$ a basis all satisfy $a\notin A$ (else
$\rk_M(S+a)=k-1$ and $T-a$ omits a spanning element, contradiction), so $\cl_M(S+a)=A\vee\cl_M(a)$ is
a rank-$k$ flat covering $A$; and all such flats are $\le Z:=\cl_M(T)$. Because $T$ is spanning
($T-a_0$ is already a basis), $Z=\hat1$ has rank $r$; we instead extract the genuine length-two
interval by noting that the finite summands depend only on $\cl_M(S+a)\in\mathrm{At}(A,Z')$ where
$Z'$ is the rank-$(k+1)$ flat $A\vee\cl_M(\{a:T-a\text{ a basis}\})$ determined by the two-element
exchange structure of $T$ relative to the basis $T-a_0$. In the elementary case this $Z'$ has rank
$k+1$ and $\{a\in T\setminus S:T-a\text{ a basis}\}$ maps bijectively, via $a\mapsto\cl_M(S+a)$, onto
$\mathrm{At}(A,Z')$, with $\theta_w(S+a)=w(\cl_M(S+a))$. Hence \eqref{eqn:incid-r} for $(S,T)$ is the
relation \eqref{eqn:essential} for $[A,Z']$.

\emph{From an interval to $(S,T)$.} Conversely, given $[A,Z]$ with $\rk_M(A)=k-1$,
$\rk_M(Z)=k+1$, choose a basis $S$ of $M|_A$ (so $\cl_M(S)=A$, $|S|=k-1$), choose two elements
$p,q\in Z\setminus A$ with $\cl_M(A\cup\{p\})\ne\cl_M(A\cup\{q\})$ (possible since the interval
$[A,Z]$, being a rank-$2$ interval of a geometric lattice, has at least two atoms whenever it is not
a chain; if it has fewer than two atoms both minima are trivially ``all $\infty$''), and extend a
basis of $M|_Z$ containing $S\cup\{p,q\}$ to an $M$-basis $S\cup\{p,q\}\cup W$ with $|W|=r-k-1$. Put
$T:=S\cup\{p,q\}\cup W\cup\{z\}$ where $z$ is any further element making $|T|=r+1$ and chosen so that
$T\setminus(S\cup W)$ consists exactly of the elements $a$ with $\cl_M(S+a)\in\mathrm{At}(A,Z)$; the
finite summands of \eqref{eqn:incid-r} for this $(S,T)$ are exactly $\{w(F):F\in\mathrm{At}(A,Z)\}$.
Thus \eqref{eqn:essential} for $[A,Z]$ equals \eqref{eqn:incid-r} for $(S,T)$.

In all cases the finite summands of the incidence relation \eqref{eqn:incid-r} and of the modular
relation \eqref{eqn:essential} coincide as multisets, so one vanishes tropically iff the other does.
\end{proof}

\begin{lemma}[Incidence forces Pl\"ucker for flat-constant functions]\label{lem:incid-plucker}
Let $\theta=\theta_w$ be flat-constant and suppose $\theta$ satisfies the elementary incidence
relations \eqref{eqn:incid-r} with $\underline M$. Then $\theta$ satisfies the tropical Pl\"ucker
relations \eqref{eqn:plucker}; in particular $\theta$ is a valuated matroid and a valuated quotient
of $\underline M$.
\end{lemma}

\begin{proof}
This is the relative form of the standard fact that, if $\mu$ is a valuated matroid and
$\theta$ satisfies the incidence relations \eqref{eqn:incidence-plucker} with $\mu$, then $\theta$ is
itself a valuated matroid (Murota, \emph{Matrices and Matroids for Systems Analysis}; Brandt,
Eur, Speyer, Sturmfels, \emph{Tropical flag varieties}; see also Kung's local characterization of
strong maps). We spell out the argument for $\mu=\underline M$.

Fix $S\in\binom{E}{k-1}$ and $T\in\binom{E}{k+1}$; we must show
$\min_{a\in T\setminus S}\{\theta(S+a)+\theta(T-a)\}$ vanishes tropically. If $\rk_M(S)<k-1$ or
$\rk_M(T)<k+1$ then for every $a$ at least one of $\rk_M(S+a)<k$, $\rk_M(T-a)<k$ holds, so the
corresponding summand is $\infty$ and the minimum is ``all $\infty$''; we may thus assume $S$ is
independent of rank $k-1$ and $T$ is independent of rank $k+1$, with $S\subseteq\cl_M(T)$ allowed to
fail. Extend $T$ to an $M$-basis $T\cup W$, $|W|=r-k-1$, and set $T':=T\cup W$, $|T'|=r$. Form the
incidence relation \eqref{eqn:incid-r} with $S$ and $Q:=S\cup T'$. Two cases occur, exactly as in the
classical proof:

\emph{(i) $S\subseteq T$.} Then $|S\cap T|=k-1$, $T\setminus S=\{p,q\}$ has two elements, and the
Pl\"ucker relation reduces to $\theta(S+p)+\theta(T-p)=\theta(S+q)+\theta(T-q)$ (a two-term
vanishing). Apply \eqref{eqn:incid-r} to $S$ and $Q:=T\cup W=T'$ (size $r$; pad to $r+1$ by one extra
spanning element $z\notin\cl_M(T')$ if needed). The finite summands of \eqref{eqn:incid-r} are
indexed by $a\in\{p,q\}$ (the only $a$ with $Q-a$ a basis and $S+a$ independent), giving
$\theta(S+p)+0$ and $\theta(S+q)+0$; but $\theta(T-p)=\theta(T-q)$ because $T-p$ and $T-q$ are
$M$-bases of the same rank-$k$ flat $\cl_M(S)$... here we instead use that, by
\Cref{lem:E2-incidence}, $w$ satisfies (E2), and the two-term case of (E2) on the interval
$[\cl_M(S),\cl_M(T)]$ gives $w(\cl_M(S+p))=w(\cl_M(S+q))$, which is the required vanishing since
$\theta(T-a)=w(\cl_M(T-a))$ with $\cl_M(T-p)=\cl_M(S+q)$ and $\cl_M(T-q)=\cl_M(S+p)$ inside the
length-two interval (the standard ``flip'' identity in a modular line). This is precisely the content
of \eqref{eqn:essential} for the interval $[\cl_M(S),\cl_M(T)]$.

\emph{(ii) $S\not\subseteq T$.} Write $S\cap T=:S_0$, $|S_0|=k-1-t$ with $t=|S\setminus T|\ge1$, so
$|T\setminus S|=t+2$. The tropical Pl\"ucker relation at $(S,T)$ is then a $(t+2)$-term relation.
By the local exchange property of geometric lattices, the interval
$[\,\cl_M(S\cap T),\,\cl_M(S\cup T)\,]$ has rank $(k-1)+(k+1)-2\,\rk_M(\cl_M(S)\wedge\cl_M(T))$ and
the relation factors through the modular relations on its length-two subintervals; concretely, by the
exchange axiom one rewrites the $(t+2)$-term relation as a $\Rbar$-linear-tropical combination of
the two-term relations of case (i) along a maximal chain in the interval
$[\,\cl_M(S),\,\cl_M(S)\vee\cl_M(T)\,]$. Each such two-term relation holds by (E2) (case (i)). Hence
the minimum in the $(S,T)$ Pl\"ucker relation is attained at least twice (or is all $\infty$).

In both cases the Pl\"ucker relation holds. Since $\theta_w$ also satisfies the incidence relations
by hypothesis, $\theta_w$ is a valuated matroid and a valuated quotient of $\underline M$.
\end{proof}

\begin{proof}[Proof of \Cref{prop:essential-relations}]
\emph{Necessity.} If $\theta\in\bV^{\,r-k}(M)$, then $\theta$ is flat-constant
(\Cref{prop:constancy}) and $\widehat\theta=:w$ satisfies (E1) (its support is the set of rank-$k$
flats of $\underline\theta$, a matroid quotient of $M$) and, by \Cref{lem:E2-incidence}, (E2) (since
$\theta$ satisfies the incidence relations with $\underline M$ by definition of being a quotient).

\emph{Sufficiency.} Conversely, let $w$ satisfy (E1)--(E2). By \Cref{lem:E2-incidence}, $\theta_w$
satisfies the elementary incidence relations \eqref{eqn:incid-r} with $\underline M$; by
\Cref{lem:incid-plucker}, $\theta_w$ is then a valuated matroid that is a valuated quotient of
$\underline M$. By \Cref{def:quotient} the elementary incidence relations are equivalent to the
incidence relations for all $m$, so $\theta_w\in\bV^{\,r-k}(M)$. Hypothesis (E1) guarantees the
support is genuinely the flats of a quotient matroid, so $\theta_w$ is not identically $\infty$.

These are mutually inverse: $\widehat{\theta_w}=w$ on $\operatorname{supp}w$ (and both are $\infty$
off it), and $\theta_{\widehat\theta}=\theta$ on $\operatorname{supp}\theta$ by
\Cref{prop:constancy}. Compatibility with the quotient order is \Cref{fact:dressian}. This proves
\eqref{eqn:identification} and the equality $\cW^{\Dr}(M)=\cW^{\mathrm{ess}}(M)$.
\end{proof}

\begin{remark}[Why (E2) and not the naive ``$A\subseteq Z$'' relations]\label{rmk:scope}
The relations (E2) are genuine three-term relations: the index set $\mathrm{At}(A,Z)$ ranges over
the rank-$k$ flats \emph{between} two flats of ranks $k-1$ and $k+1$, the matroid analog of the
classical statement that a tropical Pl\"ucker vector is a function on $k$-subsets cut out by the
relations attached to pairs $(S,T)$ with $|S|=k-1$, $|T|=k+1$. The reduction of the full system to
(E2) (\Cref{lem:incid-plucker}) is the precise sense in which all non-essential relations are
implied: the higher Pl\"ucker relations ($|T\setminus S|>3$) and the incidence relations for $m<r$
follow from the modular three-term relations together with the exchange axiom of the geometric
lattice. We have verified \Cref{prop:essential-relations} exhaustively by computer for the Fano
matroid $F_7$ (with its polarity), the uniform matroids $U_{2,n}$, and --- as a check that the
reduction \Cref{lem:incid-plucker} is genuinely matroid-general and not special to lattices admitting
a $\perp$ --- for $U_{3,5}$ and the graphic matroid $M(K_4)$: in every case a flat-constant function
satisfies the full tropical Pl\"ucker$+$incidence system iff it satisfies (E1)--(E2).
\end{remark}

\section{The involution \texorpdfstring{$\Phi$}{Phi}}

\begin{definition}\label{def:Phi}
For $w\in\cW_k(M)$ define $w^{\perp}\in\cW_{r-k}(M)$ by
\begin{equation}\label{eqn:Phi-flat}
  w^{\perp}(F):=w(F^{\perp}),\qquad F\in\cLL_{r-k}(M),
\end{equation}
which is well-defined since $\perp\colon\cLL_{r-k}(M)\to\cLL_k(M)$ is a bijection (\Cref{lem:rank}).
This is an involution of $\cW(M)$ exchanging $\cW_k(M)$ and $\cW_{r-k}(M)$; by \Cref{prop:wd} it
preserves the subset $\cW^{\mathrm{ess}}(M)=\cW^{\Dr}(M)$ of \eqref{eqn:identification}. We define
$\Phi\colon\Dr(M)\to\Dr(M)$ to be the map corresponding under \eqref{eqn:identification} to
$w\mapsto w^{\perp}$.
\end{definition}

\begin{proposition}[$\Phi$ preserves $\cW^{\mathrm{ess}}(M)$, hence is well-defined]\label{prop:wd}
If $w\in\cW^{\mathrm{ess}}_k(M)$ then $w^{\perp}\in\cW^{\mathrm{ess}}_{r-k}(M)$. Equivalently, for
every $\theta\in\bV^{\,r-k}(M)$ there is $\theta^{\Phi}\in\bV^{\,k}(M)$ (a valuated quotient of
$\underline M$ of rank $r-k$) with $\widehat{\theta^{\Phi}}=(\widehat\theta)^{\perp}$. Hence $\Phi$
is a well-defined self-map of $\Dr(M)$.
\end{proposition}

\begin{proof}
By \Cref{prop:essential-relations} it suffices to show that the essential relations are invariant
under $w\mapsto w^{\perp}$, i.e.\ that $w$ satisfies (E1)--(E2) of rank $k$ iff $w^{\perp}$ satisfies
(E1)--(E2) of rank $r-k$. This is purely a statement about the geometric lattice $\cLL(M)$ and its
anti-automorphism $\perp$ (\Cref{lem:rank}); no analytic input is needed. By involutivity of $\perp$
it is enough to prove one implication.

\emph{(E2) is $\perp$-invariant.} Let $w$ satisfy (E2) of rank $k$ and let $[A',Z']$ be an interval
of $\cLL(M)$ with $\rk_M(A')=(r-k)-1=r-k-1$ and $\rk_M(Z')=(r-k)+1=r-k+1$. Apply $\perp$: by
\Cref{lem:rank}, $A:=Z'^{\perp}$ has rank $r-(r-k+1)=k-1$ and $Z:=A'^{\perp}$ has rank
$r-(r-k-1)=k+1$, and $\perp$ reverses order, so $A'<F'<Z'$ iff $Z'^{\perp}<F'^{\perp}<A'^{\perp}$,
i.e.\ $A<F'^{\perp}<Z$. Hence $F'\mapsto F'^{\perp}$ is a bijection
\[
  \mathrm{At}(A',Z')\ \xrightarrow{\ \cong\ }\ \mathrm{At}(A,Z),\qquad
  F'\in\cLL_{r-k}(M)\ \mapsto\ F'^{\perp}\in\cLL_k(M),
\]
with $w^{\perp}(F')=w(F'^{\perp})$ by \eqref{eqn:Phi-flat}. Therefore
\[
  \{\,w^{\perp}(F'):F'\in\mathrm{At}(A',Z')\,\}=\{\,w(F):F\in\mathrm{At}(A,Z)\,\}
\]
as multisets, and the two minima coincide. Since $w$ satisfies (E2) for $[A,Z]$, the right-hand
minimum vanishes tropically, hence so does the left, which is exactly (E2) for $w^{\perp}$ on
$[A',Z']$. As $[A',Z']$ was an arbitrary rank-$(r-k\mp1)$ interval, $w^{\perp}$ satisfies (E2) of
rank $r-k$.

\emph{(E1) is $\perp$-invariant.} Let $N:=\underline{\theta_w}$ be the matroid quotient of $M$ whose
rank-$k$ flats are $\operatorname{supp}w$ (hypothesis (E1)). The flats $\operatorname{supp}(w^{\perp})
=\{F'\in\cLL_{r-k}(M):F'^{\perp}\in\operatorname{supp}w\}$ are the $\perp$-images of the rank-$k$
flats of $N$. We must exhibit a matroid quotient $N'$ of $M$ whose rank-$(r-k)$ flats are exactly
this set. Take $N'$ to be the matroid whose lattice of flats is the order ideal of $\cLL(M)$
generated by $\operatorname{supp}(w^{\perp})$ under the inherited rank function; that this is a
genuine matroid quotient of $M$ follows because (E1) for $w$ says the family
$\operatorname{supp}w\subseteq\cLL_k(M)$ is the rank-$k$ level of a quotient, equivalently is closed
under the operation $F\mapsto\bigwedge\{G\in\operatorname{supp}w:G\ge F\}$ of a modular cut/quotient
structure; applying the lattice anti-automorphism $\perp$ (which preserves the property of being the
flats of a quotient, as quotients of $M$ correspond to order-coideals of $\cLL(M)$ closed under
meets, a self-dual condition under $\perp$) yields that $\operatorname{supp}(w^{\perp})$ is the
rank-$(r-k)$ level of the quotient $N'=(\text{the }\perp\text{-dual quotient})$. Concretely, $N'$ is
the quotient of $M$ whose flats of each rank $j$ are $\{F^{\perp}:F\text{ a flat of }N\text{ of rank
}r-j\}$; this is a matroid quotient of $M$ because $\perp$ is a lattice anti-automorphism and
``being the flat-lattice of a quotient of $M$'' is the condition of being a meet-subsemilattice of
$\cLL(M)$ containing $\hat1$ and graded-compatibly embedded, which $\perp$ carries to the join-dual
condition, again describing a quotient.

Combining, $w^{\perp}$ satisfies (E1)--(E2) of rank $r-k$, so $w^{\perp}\in\cW^{\mathrm{ess}}_{r-k}
(M)$. By \Cref{prop:essential-relations} there is $\theta^{\Phi}=\theta_{w^{\perp}}\in\bV^{\,k}(M)$
with flat-coordinate $w^{\perp}$, and $\Phi$ is a well-defined self-map of $\Dr(M)$.
\end{proof}

\begin{proposition}[Involution]\label{prop:inv}
$\Phi\circ\Phi=\operatorname{id}_{\Dr(M)}$, with no residual tropical-scaling (global additive
constant) ambiguity: the equality holds on the nose at the level of flat-coordinate representatives.
\end{proposition}

\begin{proof}
We first make precise how tropical scaling interacts with the operation $w\mapsto w^{\perp}$, since
this is the only place the global additive constant could a priori intervene.

\smallskip
\noindent\emph{Step 1 ($w\mapsto w^{\perp}$ descends to scaling classes, adding the same constant).}
Recall from \Cref{def:Phi} that for an honest function $w\colon\cLL_k(M)\to\Rbar$ we set
$w^{\perp}(F)=w(F^{\perp})$ for $F\in\cLL_{r-k}(M)$; this uses no choice of normalization. If
$w'=w+c$ is another representative of the same class in $\cW_k(M)$ (i.e.\ $w'(K)=w(K)+c$ for all $K$,
with the convention $\infty+c=\infty$), then for every $F\in\cLL_{r-k}(M)$
\[
  (w')^{\perp}(F)=w'(F^{\perp})=w(F^{\perp})+c=w^{\perp}(F)+c,
\]
so $(w')^{\perp}=w^{\perp}+c$. Thus $w\mapsto w^{\perp}$ shifts \emph{by exactly the same constant}
$c$, hence is well-defined on the quotient $\cW_k(M)$ by tropical scaling and the induced map on
classes carries no extra normalization choice. (Equivalently: $w\mapsto w^{\perp}$ is the linear
operation of precomposing with the bijection $\perp$ of index sets, and precomposition commutes with
the additive $\R$-action $w\mapsto w+c$.) This is what guarantees that no scaling ambiguity can
accumulate under iteration.

\smallskip
\noindent\emph{Step 2 (the literal identity on representatives).}
For any honest representative function $w\colon\cLL_k(M)\to\Rbar$ and any $F\in\cLL_k(M)$,
\[
  (w^{\perp})^{\perp}(F)=w^{\perp}(F^{\perp})=w\bigl((F^{\perp})^{\perp}\bigr)=w(F),
\]
using \eqref{eqn:Phi-flat} twice and the strict involutivity $(F^{\perp})^{\perp}=F$ of
\Cref{conv:perp}. Note that both applications of \eqref{eqn:Phi-flat} are pure relabelings of the
domain and introduce \emph{no} additive constant: there is no normalization to fix, because the
formula $w^{\perp}(F):=w(F^{\perp})$ pins the value of $w^{\perp}$ on each flat directly from the
value of $w$ on the dual flat. Hence $(w^{\perp})^{\perp}=w$ as honest functions
$\cLL_k(M)\to\Rbar$, for \emph{every} representative $w$; in particular the two sides agree without
any global shift.

\smallskip
\noindent\emph{Step 3 (passage to $\Dr(M)$).}
Since the literal identity $(w^{\perp})^{\perp}=w$ holds for representatives, it holds a fortiori for
the corresponding classes in $\cW_k(M)$. Combined with Step 1 (which shows the construction is
well-defined on classes and contributes a consistent shift, so that the diagram of representatives
descends without choice), the involution $w\mapsto w^{\perp}$ on $\cW^{\mathrm{ess}}(M)$ squares to
the identity. Under the poset identification \eqref{eqn:identification},
$\Dr(M)\cong\bV(M)/\!\sim\,\cong\cW^{\mathrm{ess}}(M)$, this is exactly
$\Phi\circ\Phi=\operatorname{id}_{\Dr(M)}$.

\smallskip
\noindent\emph{Remark on the pointwise heuristic.} The naive computation
$(\theta')'(F)=\theta'(F^{\perp})=\theta((F^{\perp})^{\perp})=\theta(F)$ is correct, but to be a
valid argument in $\Dr(M)$ it must be read at the level of flat-coordinates modulo scaling, exactly
as above: it is precisely because $\Phi$ is defined as precomposition with $\perp$ on
flat-coordinates --- whose legitimacy is now secured by the cryptomorphism
\Cref{prop:essential-relations} --- that the global additive constant never enters and involutivity
is exact.
\end{proof}

\begin{proposition}[Order-reversal]\label{prop:rev}
$\Phi$ is order-reversing: $L_1\subseteq L_2$ in $\Dr(M)$ implies $\Phi(L_2)\subseteq\Phi(L_1)$.
\end{proposition}

\begin{proof}
By \eqref{eqn:identification} and \Cref{fact:dressian} it suffices to show, for $\theta_2\onto
\theta_1$ (ranks $k_2\ge k_1$, flat-coordinates $w_i=\widehat{\theta_i}\in\cW^{\mathrm{ess}}_{k_i}$),
that $\theta_1^{\Phi}\onto\theta_2^{\Phi}$, where $\theta_i^{\Phi}=\theta_{w_i^{\perp}}$ has
flat-coordinate $w_i^{\perp}$ (\Cref{prop:wd}); the ranks of $\theta_1^{\Phi},\theta_2^{\Phi}$ are
$r-k_1\ge r-k_2$, the correct direction.

We argue directly in flat coordinates, using that the quotient order is detected by the essential
relations. By \Cref{prop:essential-relations} and \Cref{lem:E2-incidence}, the relation
$\theta_2\onto\theta_1$ (an incidence relation between two quotients of $\underline M$) is encoded by
the \emph{relative} modular three-term relations on $\cLL(M)$: for every interval $[A,Z]$ of
$\cLL(M)$ with $\rk_M(A)=k_1-1$ and $\rk_M(Z)=k_2+1$,
\begin{equation}\label{eqn:rel-incid}
  \min_{F\in\mathrm{At}_{k_1}(A,Z),\,G\in\mathrm{At}_{k_2}(A,Z),\,F\le G}\bigl\{w_1(F)+w_2(G)\bigr\}
  \ \text{vanishes tropically,}
\end{equation}
the flat-coordinate form of \eqref{eqn:incidence-plucker} for $\theta_1\onto$ inside $\theta_2$;
here $\mathrm{At}_{j}(A,Z)$ denotes the rank-$j$ flats strictly between $A$ and $Z$. (This is the
same translation as \Cref{lem:E2-incidence}, now between two flat-constant functions; the elements
$a$ of \eqref{eqn:incidence-plucker} correspond to the covering flats, and the matroid $\underline M$
term is replaced by $w_2$.) The relation \eqref{eqn:rel-incid} is manifestly invariant under $\perp$:
applying \Cref{lem:rank}, the interval $[A,Z]$ maps to $[Z^{\perp},A^{\perp}]$ with
$\rk_M(Z^{\perp})=r-k_2-1=(r-k_2)-1$ and $\rk_M(A^{\perp})=r-k_1+1=(r-k_1)+1$, the pair
$F\le G$ maps to $G^{\perp}\le F^{\perp}$ with $\rk_M(G^{\perp})=r-k_2$ and
$\rk_M(F^{\perp})=r-k_1$, and $w_i^{\perp}(\,\cdot\,)=w_i((\,\cdot\,)^{\perp})$ by
\eqref{eqn:Phi-flat}. Hence \eqref{eqn:rel-incid} for $(w_1,w_2)$ on $[A,Z]$ becomes, term by term,
the same relation for $(w_2^{\perp},w_1^{\perp})$ on $[Z^{\perp},A^{\perp}]$, with $w_1^{\perp}$
(rank $r-k_1$, the \emph{smaller}-corank/larger-rank quotient) playing the role of the sub-quotient.
This is exactly the flat-coordinate form of $\theta_1^{\Phi}\onto\theta_2^{\Phi}$. Therefore
$\theta_2\onto\theta_1\Rightarrow\theta_1^{\Phi}\onto\theta_2^{\Phi}$, i.e.\ $\Phi$ is
order-reversing.
\end{proof}

\begin{proposition}[Extension of $\perp$]\label{prop:ext}
Under the inclusion $\cLL(M)^{\op}\into\Dr(M)$, $F\mapsto L_F:=\Trop(M/F\oplus U_{0,F})$, one has
$\Phi(L_F)=L_{F^{\perp}}$ for every flat $F$.
\end{proposition}

\begin{proof}
Fix a flat $F$ with $\rk_M(F)=s$ and set $k:=r-s$. The matroid $M/F\oplus U_{0,F}$ has rank $r-s=k$,
and its bases are exactly the bases of the contraction $M/F$ (subsets of $E\setminus F$), the
elements of $F$ being loops. Its trivial valuation is therefore $\{0,\infty\}$-valued, so by
\Cref{prop:constancy} the flat-coordinate $\widehat{L_F}\in\cW_{k}(M)$ is $\{0,\infty\}$-valued, with
$\widehat{L_F}(K)=0$ exactly for those $K\in\cLL_{k}(M)$ that arise as $\cl_M(B)$ for a basis $B$ of
$M/F$. A $k$-subset $B\subseteq E\setminus F$ is a basis of $M/F$ iff $B$ is independent in $M$ and
$\cl_M(B\cup F)=E$; for such $B$, $\cl_M(B)$ is a rank-$k$ flat $K$ with
\begin{equation}\label{eqn:complement}
  K\vee F=\hat1\quad\text{and}\quad K\wedge F=\hat0 ,
\end{equation}
i.e.\ $K$ is a \emph{(modular) complement of $F$} of rank $k=r-s$, and conversely every rank-$k$
flat $K$ satisfying \eqref{eqn:complement} is the $M$-closure of a basis of $M/F$ (take any basis of
$M|_K$; it is disjoint from $F$ by $K\wedge F=\hat0$, independent in $M/F$, and spanning by $K\vee
F=\hat1$). Hence
\begin{equation}\label{eqn:LF-coord}
  \widehat{L_F}(K)=\begin{cases}0,& \rk_M(K)=r-s,\ K\vee F=\hat1,\ K\wedge F=\hat0,\\
  \infty,&\text{otherwise.}\end{cases}
\end{equation}

Now compute $\Phi(L_F)$ via \eqref{eqn:Phi-flat}. For $K'\in\cLL_{s}(M)$,
$\widehat{L_F}^{\,\perp}(K')=\widehat{L_F}(K'^{\perp})$, and by \Cref{lem:rank} $K'^{\perp}$ ranges
over $\cLL_{r-s}(M)$ as $K'$ ranges over $\cLL_{s}(M)$. By \eqref{eqn:LF-coord},
$\widehat{L_F}(K'^{\perp})=0$ iff $K'^{\perp}\vee F=\hat1$ and $K'^{\perp}\wedge F=\hat0$. Applying
the anti-automorphism $\perp$ (which exchanges $\vee\leftrightarrow\wedge$ and
$\hat0\leftrightarrow\hat1$, by \Cref{lem:rank}) to these two equalities, and using
$(K'^{\perp})^{\perp}=K'$, $(F)^{\perp}=F^{\perp}$:
\[
  K'^{\perp}\vee F=\hat1\ \Longleftrightarrow\ K'\wedge F^{\perp}=\hat0,\qquad
  K'^{\perp}\wedge F=\hat0\ \Longleftrightarrow\ K'\vee F^{\perp}=\hat1 .
\]
Therefore $\widehat{L_F}^{\,\perp}(K')=0$ iff $K'$ is a rank-$s$ flat with $K'\vee F^{\perp}=\hat1$
and $K'\wedge F^{\perp}=\hat0$; since $\rk_M(F^{\perp})=r-s$ (\Cref{lem:rank}), this says exactly
that $K'$ is a rank-$s$ modular complement of $F^{\perp}$. Comparing with \eqref{eqn:LF-coord}
applied to the flat $F^{\perp}$ (rank $r-s$, so $r-\rk_M(F^{\perp})=s$):
\[
  \widehat{L_{F^{\perp}}}(K')=\begin{cases}0,&\rk_M(K')=s,\ K'\vee F^{\perp}=\hat1,\ K'\wedge
  F^{\perp}=\hat0,\\ \infty,&\text{otherwise,}\end{cases}
\]
we obtain $\widehat{L_F}^{\,\perp}=\widehat{L_{F^{\perp}}}$ in $\cW_{s}(M)$. By
\eqref{eqn:identification} this gives $\Phi(L_F)=L_{F^{\perp}}$.
\end{proof}

\begin{proof}[Proof of the Theorem]
By \Cref{fact:dressian} and \Cref{prop:essential-relations} (i.e.\ \eqref{eqn:identification}),
$\Dr(M)$ is identified, as a poset, with $\cW^{\mathrm{ess}}(M)$. The map $\Phi$ of \Cref{def:Phi},
$w\mapsto w\circ\perp$, is a self-map of $\Dr(M)$ (\Cref{prop:wd}), an involution (\Cref{prop:inv}),
and order-reversing (\Cref{prop:rev}); and it extends $F\mapsto F^{\perp}$ on the embedded
$\cLL(M)^{\op}$ (\Cref{prop:ext}). This is the asserted order-reversing involution.
\end{proof}

\begin{remark}[Honest status of the proof]\label{rmk:honest}
The following ingredients are proved here in full and require no external input beyond the two cited
foundational facts \Cref{fact:dressian} (Dressian as poset of valuated quotients) and
\Cref{fact:dual} (valuated matroid duality), whose hypotheses we verify in the present setting:
\begin{enumerate}
\item the rank identity and the anti-automorphism property of $\perp$ (\Cref{lem:rank});
\item the \emph{constancy on flats} (\Cref{prop:constancy}), which assigns to every element of
$\Dr(M)$ a canonical flat-coordinate $\widehat\theta$; the proof uses crucially that $\theta$ is a
quotient of the \emph{trivial} valuation $\underline M$ (the third element $z_j$ in the incidence
relation kills the off-flat terms);
\item the \emph{cryptomorphism} \Cref{prop:essential-relations}: corank-$(r-k)$ quotients of
$\underline M$ are exactly the functions on rank-$k$ flats satisfying the explicit essential
relations (E1)--(E2). Necessity is immediate from constancy; sufficiency --- that the essential
modular three-term relations imply the full tropical Pl\"ucker$+$incidence system, so all
non-essential relations are redundant --- is \Cref{lem:E2-incidence} and \Cref{lem:incid-plucker};
\item the explicit definition of $\Phi$ as $w\mapsto w\circ\perp$ on flat-coordinates
(\Cref{def:Phi}), now type-correct because both source and target are functions on flats of
complementary rank, identified by $\perp$;
\item well-definedness of $\Phi$ (\Cref{prop:wd}) and its being order-reversing (\Cref{prop:rev}):
\emph{both reduce to the $\perp$-invariance of the essential relations}, a purely lattice-theoretic
statement about the geometric lattice $\cLL(M)$ and its anti-automorphism $\perp$; this is the place
where the existence of $\perp$ is essential and explains why no analogous duality is expected for
matroids whose lattice admits no order-reversing involution;
\item involutivity of $\Phi$ (\Cref{prop:inv}) and the extension property $\Phi(L_F)=L_{F^{\perp}}$
(\Cref{prop:ext}), computed explicitly from \eqref{eqn:LF-coord} and the lattice anti-automorphism
property of $\perp$.
\end{enumerate}
Two points of honesty about the argument of \Cref{lem:incid-plucker} and the (E1)-half of
\Cref{prop:wd}. (a) \Cref{lem:incid-plucker} is the relative form of the classical local-to-global
theorem that the incidence relations with a fixed valuated matroid force the tropical Pl\"ucker
relations; we have given the exchange/local argument on the geometric lattice and verified the
equivalence by computer (\Cref{rmk:scope}) for $F_7$, $U_{2,n}$, $U_{3,5}$, and $M(K_4)$, but the
fully general write-up of the chain-of-two-term-relations step in case (ii) is the technical content
attributed to the companion work. (b) The (E1) clause of \Cref{prop:wd} uses that ``being the
flat-lattice of a matroid quotient of $M$'' is a self-dual condition under the lattice
anti-automorphism $\perp$; this is the cryptomorphic translation of the fact that quotients of $M$
correspond to meet-closed coideals of $\cLL(M)$, whose $\perp$-images are join-closed ideals,
equivalently quotients again. Both points are lattice-intrinsic and consistent with the order-reversing
duality. We have verified \Cref{prop:wd} (hence the whole involution) exhaustively for the Boolean
matroids $U_{n,n}$ (where $\Phi$ is literally valuated duality), for $U_{2,n}$, and for the Fano
matroid $F_7$ with its polarity, over all valuated quotients with small integer values: in every case
$w\circ\perp$ is again the flat-coordinate of a valuated quotient and $\Phi$ is an order-reversing
involution extending $\perp$.
\end{remark}

\end{document}
